\subsection{A joint amplitude--phase consistency relation}
\label{sec:kernel-consistency}

The lensing response constrains amplitude and phase jointly.  In
particular, a power-spectrum-independent variance relation was established
for weak lensing by Inamori and Suyama~\cite{inamori2021} and extended to
weak perturbations of a fixed strong lens by Nakazono and
Suyama~\cite{nakazono2026}.  In the single-plane Born model used here,
the underlying identity follows directly from the kernel of
Oguri and Takahashi~\cite{oguri2022}:
\begin{equation}
 H_A(f,k,\theta)=b(\theta)
 \frac{\sin s+i(\cos s-1)}{s},\qquad
 s=\frac{r_F^2(f)k^2b(\theta)}{2}.
\end{equation}
Since $r_F^2(2f)=r_F^2(f)/2$, elementary trigonometry gives
\begin{equation}
 |H_A(f,k,\theta)|^2
 =\frac{4b(\theta)^2\sin^2(s/2)}{s^2}
 =\left[\operatorname{Re}H_A(2f,k,\theta)\right]^2.
 \label{eq:modewise-consistency}
\end{equation}
The result includes saddles; $b$ need not be positive.  Integrating
against any nonnegative homogeneous convergence power spectrum yields,
to leading nonzero order in the perturbation,
\begin{equation}
 \left\langle K(f)^2\right\rangle+
 \left\langle S(f)^2\right\rangle
 =\left\langle K(2f)^2\right\rangle,
 \label{eq:ensemble-consistency}
\end{equation}
where $K=\operatorname{Re}\eta$ and $S=\operatorname{Im}\eta$ at first
order, after removing the geometric arrival-time phase.  More generally,
the quantities in Ref.~\cite{nakazono2026} are defined through the
logarithm of the amplification factor; their leading variances agree
with those used here.  The brackets average over weak perturbations at
fixed macro lens.  The code tests Eq.~\eqref{eq:modewise-consistency}
directly, including negative-parity cases.

Nakazono and Suyama~\cite{nakazono2026} also obtain, through second
order in the weak potential and with the smooth macro lens held fixed,
\begin{equation}
 \langle S(f)\rangle-\tfrac12\langle S(2f)\rangle
   =-\langle K(f)S(f)\rangle,\qquad
 \langle K(f)\rangle=-\langle K(f)^2\rangle.
\end{equation}
Here $F_j/F_{j,0}=\exp[K_j+i(S_j+2\pi f\Delta t_j)]$ defines the
logarithmic amplitude and the delay-subtracted phase.  Within their
geometric macro-lens and Limber treatment, they additionally derive
$\langle e^{2K_j}\rangle=1$ beyond the perturbative expansion.
Our first-order covariance calculation does not predict the second-order
means; testing these identities requires a consistently extended
propagation calculation and the same ensemble definition.  They cannot
be applied unchanged to projected image differences.

Three limitations matter observationally.  First, the relation holds
for ensemble moments, not the squared residual of one realization:
cross terms between different Fourier modes cancel in the ensemble,
not in each event.  Second, our detector-weighted projection $Q$ removes
amplitude, phase, and delay using a finite frequency band.  It does not
preserve Eq.~\eqref{eq:ensemble-consistency}; a valid test must pass the
theoretical joint covariance through the same projection, frequency
window, and noise weighting as the data.  Third, independent image
residuals can have their variances added, but cross-correlated images
with different macro kernels introduce cross terms that need not obey
the simple relation.  It would therefore be incorrect to apply it
unchanged to a measured, projected image difference.

These relations restrict a lensing model, but do not identify lensing
uniquely.  A multiplicative detector response can also create correlated
amplitude and phase errors.  Stronger evidence requires a common
astrophysical transfer function across independent detectors, consistency
with lens geometry and image parity, and calibration information.

\paragraph{Causality and image ratios.}
Kramers--Kronig relations constrain the full physical lens transfer
function when its time origin is fixed at the earliest arrival.
They follow from causality~\cite{tanaka2023} but cannot automatically
be imposed on the normalized ratio $F_1/F_2$:
division introduces the zeros of $F_2$ as possible poles, and the inverse
response need not be causal.  An isolated, infinite quadratic saddle
model also has no lower bound on its local Fermat time, so it does not
by itself supply the global causal transfer function.  Finite-band
projection changes the relation further.  Finally, recovering phase
from magnitude requires any zero-dependent all-pass contributions
\cite{suyama2025}, and
causal instrumental filters can obey dispersion relations too.  Thus
causality is a physical consistency requirement, not a unique detector
systematics veto.
